\documentclass[12pt]{article}
\usepackage[T1]{fontenc}
\usepackage[utf8]{inputenc}
\usepackage{lmodern}
\usepackage{textcomp}
\usepackage{enumitem}
\usepackage{geometry}
\geometry{margin=1in}
\begin{document}
\begin{center}
Answer Key - Political Science - Graduate Level Assessment 22 \
\small (Answers are ordered exactly as in the question paper.)
\end{center}

\section*{Multiple Choice}
\begin{enumerate}[label=\arabic*.]
\item \textbf{Answer:} D
\\ \textit{Options:} A) Foreign policies of other states; B) International law; C) Intergovernmental organizations; D) All of the above
\\ \textit{Note:} The correct answer is "All of the above" because US foreign policy decision-making is influenced by a multitude of factors, including economic conditions, public opinion, legal frameworks, and international relations, all of which can serve as constraints.
\item \textbf{Answer:} A
\\ \textit{Options:} A) Richard Nixon; B) George H. W. Bush; C) Jimmy Carter; D) Ronald Reagan
\\ \textit{Note:} Richard Nixon was the first American president to visit communist China in 1972, marking a significant shift in U.S.-China relations and a strategic move in the context of the Cold War.
\item \textbf{Answer:} A
\\ \textit{Options:} A) Voluntary reliance on an external power for security; B) Willful openness to colonization; C) Cultural imperialism; D) Open advocacy of imperialism for economic gain
\\ \textit{Note:} The phrase 'empire by invitation' refers to a scenario where a state voluntarily seeks the protection or governance of a more powerful external entity, thereby creating a dependent relationship for security and stability.
\item \textbf{Answer:} A
\\ \textit{Options:} A) George Kennan; B) John Foster Dulles; C) Henry Kissinger; D) Dwight Eisenhower
\\ \textit{Note:} George Kennan is considered the "father" of containment due to his influential 1946 "Long Telegram" and subsequent writings that articulated the strategy of preventing the spread of communism during the Cold War.
\item \textbf{Answer:} D
\\ \textit{Options:} A) Naval victories; B) Diplomacy; C) British preoccupation with Europe; D) All of the above
\\ \textit{Note:} The correct answer "All of the above" reflects the multifaceted nature of the factors that contributed to the US avoiding catastrophe in 1814, including effective military leadership, strategic defensive measures, and the limitations of British resources and priorities, all of which collectively mitigated the impact of British aggression during the War of 1812.
\end{enumerate}

\section*{True / False}
\begin{enumerate}[label=\arabic*.]
\setcounter{enumi}{5}
\item \textbf{Answer:} True
\\ \textit{Options:} A) True; B) False
\\ \textit{Note:} Vectors of economic incentives directly influence individual and group choices by aligning personal financial interests with political actions, whereas institutional constraints, information, and goals serve as contextual frameworks that shape behavior without inherently motivating it through economic benefit.
\item \textbf{Answer:} True
\\ \textit{Options:} A) True; B) False
\\ \textit{Note:} The statement is true because it reflects Carl von Clausewitz's assertion that war is an extension of political discourse, where states resort to military action as a means to pursue and accomplish their strategic goals when diplomatic efforts fail.
\item \textbf{Answer:} False
\\ \textit{Options:} A) True; B) False
\\ \textit{Note:} The National Security Council is primarily a coordinating body that advises the President on national security and foreign policy matters, while the actual design and execution of U.S. foreign policy initiatives are primarily the responsibilities of the State Department and other executive agencies.
\item \textbf{Answer:} False
\\ \textit{Options:} A) True; B) False
\\ \textit{Note:} The term 'Rogue States' refers to nations that are considered a threat to international peace and security due to their policies or actions, and it is not limited to communist nations or those aligned with the Soviet Union during the Cold War.
\item \textbf{Answer:} False
\\ \textit{Options:} A) True; B) False
\\ \textit{Note:} George H.W. Bush's administration actually supported Boris Yeltsin's reforms during the transition from communism in Russia, while maintaining a more cautious relationship with Mikhail Gorbachev, especially as Gorbachev's policies began to falter.
\end{enumerate}

\section*{Long Form Response}
\begin{enumerate}[label=\arabic*.]
\setcounter{enumi}{10}
\item \textbf{Answer:} The differing relationships between U.S. and European security studies can be attributed to a confluence of historical, political, and cultural factors. Historically, the U.S. emerged as a global power post-World War II, shaping its security approach around deterrence and military intervention, while Europe, with its emphasis on multilateralism and collective security, often reflects its experiences with devastating conflicts. Politically, the U.S. tends to prioritize unilateralism and a more aggressive security posture, influenced by its domestic politics and global ambitions, whereas European nations generally advocate for diplomacy and consensus-building within frameworks like NATO and the EU. Culturally, differing perceptions of security risks and the role of military force play a significant role, with European states often more cautious about military engagements due to their historical contexts. Together, these factors create a complex web that defines the distinct yet interrelated security studies landscape in the U.S. and Europe.
\\ \textit{Note:} The answer correctly identifies that the U.S. and European security studies diverge due to historical contexts of power dynamics, differing political priorities such as unilateralism versus multilateralism, and cultural experiences shaped by their unique histories of conflict and cooperation.
\item \textbf{Answer:} The concept of the 'subaltern' in postcolonialism refers to populations that are marginalized or exist outside of hegemonic power structures, highlighting their lack of agency and representation within dominant narratives. This notion underscores the ways in which colonial and postcolonial contexts have perpetuated the silencing and exclusion of these groups, often rendering their voices and experiences invisible. The implications for marginalized populations are profound, as it compels scholars and activists to critically examine the dynamics of power, identity, and resistance. By acknowledging the subaltern, one can better understand the complexities of social justice movements and the necessity for inclusive frameworks that elevate the perspectives of those historically subordinated. Ultimately, the subaltern concept challenges the prevailing discourses and prompts a reevaluation of how power operates in society.
\\ \textit{Note:} The answer correctly identifies the 'subaltern' as marginalized groups lacking agency within hegemonic power structures, emphasizing the historical silencing of their voices and experiences in colonial and postcolonial contexts, which is crucial for understanding the implications for these populations.
\end{enumerate}

\vfill
\noindent\textit{End of Answer Key}
\end{document}